\subsection{Enunciado}

\begin{enunciado}{Proceso metalmecanico}
Se toman cuatro submuestras de cuatro mediciones:
\[
(0.5,0.6,0.4,0.3),\quad
(0.5,0.5,0.4,0.6),
\]
\[
(0.7,0.5,0.5,0.7),\quad
(0.5,0.5,0.5,0.5).
\]

\begin{enumerate}[label=(\alph*)]
\item Con esta informacion encuentre los limites para el promedio y, a su vez, para la diferencia entre el maximo y minimo.

\item Si hoy toma una muestra la cual entrega los siguientes resultados:
\[
0.4\ \text{mm},\quad 0.7\ \text{mm},\quad 0.5\ \text{mm},\quad 0.9\ \text{mm}.
\]
Basandose en la grafica anteriormente realizada, ¿que puede decir del proceso?

\item Usted cree que debe controlar mejor lo que recibe como insumo, para lo cual desea proponer un sistema de muestreo. Para ello define un \(AQL\) de \(0.02\) y un \(LPTD\) de \(0.08\), con un \(\alpha\) de \(0.05\) y un \(\beta\) de \(0.1\). Indique y explique el plan de muestreo.
\end{enumerate}
\end{enunciado}
